\documentclass[12pt]{article}

% Packages
\usepackage{geometry}
\usepackage{graphicx}
\usepackage{booktabs}
\usepackage{hyperref}
\usepackage{float}
\usepackage{listings}
\usepackage{xcolor}

\geometry{margin=1in}

\title{
Hands-On MIS Projects\\
Chapter 1 - Information Systems in Global Business Today
}

\author{
Parsa Mohebian
}

\date{\today}


\begin{document}

\maketitle

\tableofcontents

\newpage


\section{Introduction}

This report analyzes management information system problems and demonstrates
how information systems improve decision making using transactional sales data.


\section{Management Decision Problem 1-7: Snyder's of Hanover}

\subsection{Problem Summary}

Snyder's financial department depended heavily on spreadsheets and manual processes
for collecting and reporting financial information. Every month, the financial analyst
collected spreadsheets from more than 50 departments worldwide, consolidated the data,
and created the company's profit-and-loss statement manually. This process was slow,
difficult to maintain, and vulnerable to errors.

\subsection{Impact on Business Performance}

The spreadsheet-based system negatively affected business performance.

\begin{itemize}
    \item The financial analyst spent significant time collecting and consolidating
    spreadsheets instead of performing financial analysis.
    
    \item Manual data entry increased the possibility of errors and inconsistent reports.
    
    \item Financial information was delayed because reports required manual preparation.
    
    \item The company faced higher operational costs because employees performed
    repetitive administrative tasks.
    
    \item The system was difficult to scale as the company expanded globally.
\end{itemize}


\subsection{Impact on Management Decision Making}

The existing system limited management decision making because managers did not have
fast access to accurate and updated financial information.

Delayed reporting made it harder to identify financial problems, compare department
performance, and respond quickly to changes in business conditions. Incorrect or
outdated information could result in poor decisions regarding budgeting, production,
and resource allocation.


\subsection{Recommended Information System Solution}

Snyder's should implement an integrated financial information system based on a
centralized database. The system should allow departments to enter financial data
directly into a shared platform and automatically generate standardized reports.

The system should include automated reporting, security controls, and management
dashboards. An Enterprise Resource Planning (ERP) system could integrate financial
data with other business functions such as sales, inventory, and production.

This solution would reduce manual work, improve reporting accuracy, provide faster
access to information, and support better management decisions.

\section{Management Decision Problem 1-8: Dollar General}

\subsection{Problem Summary}

Dollar General operates discount stores where minimizing costs is essential.
However, the company does not have an automated inventory tracking system at
each store. Managers cannot accurately verify incoming deliveries, which has
resulted in inventory errors and increasing merchandise losses.

The company needs to determine whether an information system investment can
improve inventory accuracy, reduce losses, and support better operational
decisions.

\subsection{Decisions Before Investing in an Information System}

Before investing in an inventory management system, Dollar General should
evaluate several important factors.

\begin{itemize}

\item \textbf{Identify business goals:}
The company should define the problems it wants to solve, such as reducing
inventory losses, improving inventory accuracy, and ensuring product availability.

\item \textbf{Determine required system features:}
The company should decide whether it needs barcode scanning, real-time inventory
tracking, automatic stock updates, inventory alerts, and reporting tools.

\item \textbf{Evaluate costs and benefits:}
Management should compare implementation costs, including hardware, software,
and training, against expected benefits such as lower losses and improved
operational efficiency.

\item \textbf{Plan implementation:}
The company must decide how the system will be introduced, how employees will be
trained, and how existing inventory data will be transferred.

\item \textbf{Consider security requirements:}
The company should establish rules for protecting inventory data and controlling
access to the system.

\end{itemize}

Dollar General should implement an automated inventory management system using
barcode scanning and a centralized database. This system would provide accurate
inventory information, reduce merchandise losses, improve replenishment decisions,
and increase overall operational efficiency.
\section{Database Analysis Project 1-9}

\subsection{Dataset Description}

The project uses a transactional marketplace database stored in SQLite format.
The database contains raw sales transaction data that can be used to analyze
product demand, regional sales performance, seller performance, and monthly sales
trends.

The database contains five main tables:

\begin{itemize}
    \item \texttt{orders}: contains order identification numbers, order dates,
    and customer contact information.
    
    \item \texttt{order\_items}: contains the products sold in each order,
    seller identifiers, product identifiers, and item prices.
    
    \item \texttt{products}: contains product identifiers, product categories,
    and product weight information.
    
    \item \texttt{sellers}: contains seller identifiers and seller state
    information.
    
    \item \texttt{order\_reviews}: contains customer review scores and review
    comments.
\end{itemize}

The main analysis used the \texttt{orders}, \texttt{order\_items},
\texttt{products}, and \texttt{sellers} tables. The \texttt{order\_reviews}
table was not required for the three main sales trend questions, but it could be
used in future analysis to study customer satisfaction.

The dataset included 96,478 orders, 110,197 order item records, 32,216 products,
2,970 sellers, and 39,462 reviews. The order dates ranged from September 15,
2016 to August 29, 2018.

\begin{table}[H]
\centering
\begin{tabular}{lr}
\toprule
Table & Number of Records \\
\midrule
\texttt{orders} & 96,478 \\
\texttt{order\_items} & 110,197 \\
\texttt{products} & 32,216 \\
\texttt{sellers} & 2,970 \\
\texttt{order\_reviews} & 39,462 \\
\bottomrule
\end{tabular}
\caption{Database Tables and Record Counts}
\end{table}


\subsection{Data Preparation}

The original dataset was provided as a SQLite database file. Because the project
requires database querying and reporting, SQLite was used to inspect the database
structure and run SQL queries. The database was also exported to CSV files so the
tables could be opened in spreadsheet or database software if needed.

A Python script was created to export each SQLite table into a separate CSV file.
The exported files were:

\begin{itemize}
    \item \texttt{orders.csv}
    \item \texttt{order\_items.csv}
    \item \texttt{products.csv}
    \item \texttt{sellers.csv}
    \item \texttt{order\_reviews.csv}
\end{itemize}

The main sales analysis was performed by joining the order, product, and seller
tables into a single SQL view named \texttt{sales\_analysis\_base}. This view
created a structure similar to the textbook sales database.

\begin{table}[H]
\centering
\begin{tabular}{ll}
\toprule
Textbook Field & Dataset Field Used \\
\midrule
Store identification number & \texttt{seller\_id} \\
Sales region & \texttt{seller\_state} \\
Item number & \texttt{product\_id} \\
Item description & \texttt{product\_category\_name} \\
Unit price & \texttt{price} \\
Units sold & Count of records in \texttt{order\_items} \\
Sales period & Month extracted from \texttt{timestamp} \\
\bottomrule
\end{tabular}
\caption{Mapping Textbook Sales Fields to Dataset Fields}
\end{table}

A missing data check was also performed on the joined sales view. The analysis
found no missing product categories, seller states, or prices in the final
analysis table.


\subsection{Sales Analysis Queries}

SQL queries were used to answer the three management questions in the assignment.
The queries joined the transactional tables and summarized the data by product
category, seller region, seller, and month.

The first query group identified which products should be restocked by ranking
product categories and individual products by units sold and total revenue. Since
the dataset did not include inventory quantity, high sales volume was used as a
proxy for high customer demand.

The second query group identified stores and regions that may benefit from
promotional campaigns. In this dataset, \texttt{seller\_id} represented the
store or seller, and \texttt{seller\_state} represented the sales region. Regions
with low revenue and low sales volume were treated as candidates for additional
marketing support.

The third query group analyzed monthly sales trends. Since the dataset did not
include a direct discount field, monthly revenue and units sold were used to
identify stronger and weaker sales periods. Stronger months were considered
better candidates for full-price selling, while weaker months were considered
better candidates for discounts or promotional campaigns.

\subsection{Restocking Analysis}

The first management question asks which products should be restocked. The dataset
does not include current inventory levels or stock quantities. Therefore, units sold
were used as a proxy for customer demand. Product categories with the highest sales
volume were treated as the strongest candidates for restocking.

The analysis showed that the highest-demand product categories were
\texttt{cama\_mesa\_banho}, \texttt{beleza\_saude}, \texttt{esporte\_lazer},
\texttt{moveis\_decoracao}, and \texttt{informatica\_acessorios}.

\begin{table}[H]
\centering
\begin{tabular}{lrr}
\toprule
Product Category & Units Sold & Total Revenue \\
\midrule
\texttt{cama\_mesa\_banho} & 10,953 & 1,023,434.76 \\
\texttt{beleza\_saude} & 9,465 & 1,233,131.72 \\
\texttt{esporte\_lazer} & 8,431 & 954,852.55 \\
\texttt{moveis\_decoracao} & 8,160 & 711,927.69 \\
\texttt{informatica\_acessorios} & 7,644 & 888,724.61 \\
\bottomrule
\end{tabular}
\caption{Top Product Categories by Units Sold}
\end{table}

Based on this analysis, managers should prioritize restocking these high-demand
categories. The category \texttt{beleza\_saude} is especially important because it
has both high sales volume and the highest total revenue among the top categories.
This suggests that keeping this category available could have a significant impact
on sales performance.

\subsection{Promotional Campaign Analysis}

The second management question asks which stores and sales regions would benefit
from a promotional campaign and additional marketing. In this dataset,
\texttt{seller\_id} was used as the store identifier, and \texttt{seller\_state}
was used as the sales region.

The analysis compared sales performance by region using units sold, number of
sellers, total revenue, and average price. Regions with low total revenue were
identified as possible candidates for promotional campaigns or additional
marketing support.

\begin{table}[H]
\centering
\begin{tabular}{lrrr}
\toprule
Region & Units Sold & Number of Sellers & Total Revenue \\
\midrule
AM & 3 & 1 & 1,177.00 \\
PA & 8 & 1 & 1,238.00 \\
SE & 10 & 2 & 1,606.20 \\
PI & 11 & 1 & 2,383.00 \\
RO & 14 & 2 & 4,762.20 \\
MS & 50 & 5 & 8,551.69 \\
RN & 56 & 5 & 9,992.60 \\
PB & 37 & 6 & 16,395.50 \\
MT & 144 & 4 & 17,007.82 \\
CE & 90 & 12 & 19,517.84 \\
\bottomrule
\end{tabular}
\caption{Lowest-Performing Sales Regions by Revenue}
\end{table}

The regions \texttt{AM}, \texttt{PA}, \texttt{SE}, \texttt{PI}, and
\texttt{RO} had the lowest total revenue. These regions may benefit from
additional marketing campaigns because their sales activity is much lower than
the strongest regions.

However, the results should be interpreted carefully. Some of these regions have
only one or two sellers, which suggests that low revenue may be caused by weak
market presence rather than only weak customer demand. Therefore, management
should consider both promotional campaigns and seller expansion strategies in
these regions.

For comparison, the strongest region was \texttt{SP}, with 78,604 units sold and
8,509,511.46 in total revenue. This shows a large performance gap between
\texttt{SP} and the lower-performing regions.

\subsection{Pricing and Discount Timing Analysis}

The third management question asks when products should be offered at full price
and when discounts should be used. The dataset does not contain a direct discount
field, so monthly sales volume and total revenue were used to identify stronger
and weaker sales periods.

High-demand months are better candidates for full-price selling because customers
are already purchasing more products. Low-demand months are better candidates for
discounts, promotions, or additional marketing because sales activity is weaker.

\begin{table}[H]
\centering
\begin{tabular}{lrr}
\toprule
Month & Units Sold & Total Revenue \\
\midrule
May & 11,814 & 1,466,882.94 \\
August & 11,939 & 1,393,276.34 \\
July & 11,379 & 1,349,557.98 \\
April & 10,396 & 1,314,203.77 \\
March & 10,914 & 1,312,555.10 \\
\bottomrule
\end{tabular}
\caption{Strongest Months by Total Revenue}
\end{table}

Based on this result, managers should consider selling products closer to full
price during March, April, May, July, and August. These months had the strongest
sales performance and may not require aggressive discounting.

\begin{table}[H]
\centering
\begin{tabular}{lrr}
\toprule
Month & Units Sold & Total Revenue \\
\midrule
September & 4,740 & 607,534.64 \\
October & 5,527 & 688,572.76 \\
December & 6,188 & 726,044.09 \\
November & 8,475 & 987,765.37 \\
January & 8,950 & 1,036,443.36 \\
\bottomrule
\end{tabular}
\caption{Weakest Months by Total Revenue}
\end{table}

Discounts and promotional campaigns should be considered during September,
October, December, November, and January because these months had weaker sales
performance. Promotions during these periods may help increase demand and improve
revenue.

The results should be interpreted carefully because the dataset does not cover
all calendar months equally. The data begins in September 2016 and ends in August
2018, so some periods may be incomplete. Even with this limitation, the monthly
trend provides useful guidance for pricing and promotional planning.

\section{Job Search Project 1-10}

% Added later


\section{Conclusion}

% Final conclusion


\end{document}